\section*{الفصل \ref{ch:b1:precipitation} --- الترسيب والذوبانية}

\begin{solution}{exo:b1:precipitation:1}
\ce{BaSO4(s) <=> Ba^2+ + SO4^2-}، $K_s = [\ce{Ba^2+}][\ce{SO4^2-}] =
10^{-9.97}$. \ce{CaF2(s) <=> Ca^2+ + 2F-}، $K_s = [\ce{Ca^2+}][\ce{F-}]^2 =
10^{-9.84}$. \ce{Ag2CrO4(s) <=> 2Ag+ + CrO4^2-}، $K_s =
[\ce{Ag+}]^2[\ce{CrO4^2-}] = 10^{-11.95}$. \ce{Fe(OH)3(s) <=> Fe^3+ + 3OH-}،
$K_s = [\ce{Fe^3+}][\ce{OH-}]^3 = 10^{-38.55}$. ومن أجل $K_s = \num{2.0e-13}$،
$\mathrm{p}K_s = 12.70$.
\end{solution}

\begin{solution}{exo:b1:precipitation:2}
$s = \sqrt{K_s} = 10^{-12.27/2} = \qty{7.3e-7}{mol/L}$، أي
$\num{7.3e-7} \times 187.8 = \qty{1.4e-4}{g/L} = \qty{0.14}{mg/L}$.
\end{solution}

\begin{solution}{exo:b1:precipitation:3}
بعد الخلط (الحجم مضاعف)، $[\ce{Ba^2+}] = \qty{1.0e-3}{mol/L}$
و$[\ce{SO4^2-}] = \qty{1.0e-4}{mol/L}$: $Q_s = \num{1.0e-7} > K_s =
\num{1.1e-10}$، فتترسب كبريتات الباريوم. وفي الحالة الثانية
$[\ce{SO4^2-}] = \num{1.0e-7}$ و$Q_s = \num{1.0e-10} < K_s$: لا
راسب، وإن كان ذلك بفارق ضئيل.
\end{solution}

\begin{solution}{exo:b1:precipitation:4}
في الماء $s = 10^{-9.97/2} = \qty{1.0e-5}{mol/L}$. وفي كبريتات تركيزها \qty{0.010}{mol/L}
(\cref{prop:b1:precipitation:common-ion})، $s = 10^{-9.97}/0.010 =
\qty{1.1e-8}{mol/L}$، أي أقل بنحو ألف مرة (العامل 970).
\end{solution}

\begin{solution}{exo:b1:precipitation:5}
في الماء $s = (10^{-9.84}/4)^{1/3} = \qty{3.3e-4}{mol/L}$، أي
\qty{26}{mg/L}. وفي كلوريد الكالسيوم تركيزه \qty{0.10}{mol/L}،
$[\ce{Ca^2+}] \approx 0.10$ و$[\ce{F-}] = 2s$، فيكون
$s = \frac12\sqrt{10^{-9.84}/0.10} = \qty{1.9e-5}{mol/L}$.
والتقريب $s \ll 0.10$ صحيح بفارق أربع رتب قدر.
\end{solution}

\begin{solution}{exo:b1:precipitation:6}
البدء: $\mathrm{pH} = 14.00 - \frac12(11.25 + \log 0.050) = 14.00 - 4.97 =
9.03$. وعند \qty{99}{\percent}، $[\ce{Mg^2+}] = \num{5.0e-4}$
و$\mathrm{pH} = 14.00 - \frac12(11.25 - 3.30) = 10.03$. وعلى محور pH،
يوجد \ce{Mg(OH)2} فوق 9.03؛ وتحتها يكون المغنيسيوم كله \ce{Mg^2+}
بتركيز \qty{0.050}{mol/L}.
\end{solution}

\begin{solution}{exo:b1:precipitation:7}
يظهر بروميد الفضة أولًا، عند $[\ce{Ag+}] = 10^{-12.27 + 2} =
10^{-10.27}$، وكلوريد الفضة عند $10^{-7.75}$. وحين يظهر كلوريد الفضة
يكون $[\ce{Br-}] = 10^{-12.27}/10^{-7.75} = 10^{-4.52} =
\qty{3.0e-5}{mol/L}$، أي \qty{0.30}{\percent} من البروميد. ووفق المعيار
المعتاد \qty{0.1}{\percent} لا يكون الفصل مقبولًا: فقيمتا
$\mathrm{p}K_s$ لا تختلفان إلا بمقدار 2.5.
\end{solution}

\begin{solution}{exo:b1:precipitation:8}
تتكوّن كرومات الفضة حين $[\ce{Ag+}]^2 \times \num{5.0e-3} = 10^{-11.95}$،
$[\ce{Ag+}] = \qty{1.5e-5}{mol/L}$. ومن ثم $[\ce{Cl-}] = 10^{-9.75}/\num{1.5e-5}
= \qty{1.2e-5}{mol/L}$، أي كسر \num{2.4e-4} (\qty{0.024}{\percent}) من
الكلوريد: فيظهر اللون الأحمر حين يكون الكلوريد قد ترسب
كله عمليًا.
\end{solution}

\begin{solution}{exo:b1:precipitation:9}
$[\ce{CO3^2-}] = [\ce{HCO3-}]K_{a2}/[\ce{H3O+}] = [\ce{HCO3-}]\,10^{\mathrm{pH}
- \mathrm{p}K_{a2}}$، فيكون $K_s = [\ce{Ca^2+}][\ce{HCO3-}]\,10^{\mathrm{pH} -
\mathrm{p}K_{a2}}$، وهي العلاقة المطلوبة. ومع $[\ce{Ca^2+}] =
[\ce{HCO3-}] = s$: $s^2 = 10^{-8.30 + 10.33 - 8.30} = 10^{-6.27}$،
$s = \qty{7.3e-4}{mol/L}$ (\qty{73}{mg/L} من \ce{CaCO3})، مقابل
$10^{-4.15} = \qty{7.1e-5}{mol/L}$ لو لم يكن أيون الكربونات قاعدة.
\end{solution}

\begin{solution}{exo:b1:precipitation:10}
1.~بإضافة \ce{CO2(g) <=> CO2(aq)}: $K' = K K_H = 10^{-4.34 - 1.47} =
10^{-5.81}$.
2.~$[\ce{Ca^2+}] = s$ و$[\ce{HCO3-}] = 2s$، فيكون $4s^3 = K'p$. وعند
\qty{0.010}{bar}: $s = (10^{-5.81} \times 0.010/4)^{1/3} =
\qty{1.6e-3}{mol/L}$، أي \qty{157}{mg/L}. وتحت الهواء الطلق:
$s = \qty{5.5e-4}{mol/L}$، أي \qty{55}{mg/L}. (تحقق: تحت هواء التربة
$[\ce{CO2(aq)}] = 10^{-3.47}$ و$\mathrm{pH} = 6.37 + \log(3.14\times
10^{-3}/3.39 \times 10^{-4}) = 7.34$: فالهيدروجينوكربونات هي الغالبة فعلًا.)
3.~يبلغ الماء المشبع بالكالسيت تحت ضغط ثنائي أكسيد الكربون العالي في
التربة سقفَ المغارة، حيث الضغط أدنى: فيفقد
ثنائي أكسيد الكربون، ويصير $Q > K'$، فيترسب الكالسيت حيث تتدلى القطرة،
حلقة بعد حلقة.
% ledger: air:CO2, dfg:CO2_g, dfg:CO2_ao (log KH computed in figdata/bachelor-1/aqueous-constants.py)
\end{solution}

\begin{solution}{exo:b1:precipitation:11}
1.~البدء عند $\mathrm{pH} = 14.00 - \frac12(16.38 - 2.00) = 6.81$. وإعادة
الذوبان التامة حين $[\ce{Zn(OH)4^2-}] = 10^{-2} = 10^{-1.93}[\ce{OH-}]^2$،
$[\ce{OH-}] = 10^{-0.035}$، $\mathrm{pH} = 13.97$. فالهيدروكسيد موجود
بين pH 6.81 و13.97.
2.~$\mathrm{pH} = 14.00 - 14.45/4 = 10.39$، $s_{\min} = 2 \times
10^{-(16.38 + 1.93)/2} = \qty{1.4e-9}{mol/L}$ (في نموذج النوعين هذا).
3.~مستقيم ميله $-2$ من $(6.81, -2)$ نزولًا إلى القيمة الصغرى قرب
$(10.39, -8.9)$، ثم ميل $+2$ صعودًا إلى $(13.97, -2)$؛ أما المنحنى الحقيقي
فتدوّره الأنواع الهيدروكسيلية الأخرى
(الشكل في \cref{sec:b1:precipitation:amphoteric}).
\end{solution}

\begin{solution}{exo:b1:precipitation:12}
1.~\ce{Al(OH)3(s) + OH- <=> Al(OH)4-}، $K = 10^{-1.23}$؛ ويكون كل
الألومنيوم مذابًا حين $10^{-3} \leqslant K[\ce{OH-}]$، أي
$[\ce{OH-}] \geqslant 10^{-1.77}$، $\mathrm{pH} \geqslant 12.23$.
2.~$[\ce{Fe^3+}] = 10^{-38.55}/10^{-3 \times 1.77} = 10^{-33.2}$: يبقى الحديد
كله على هيئة \ce{Fe(OH)3}؛ ويفصل الترشيح الحديد (الجسم الصلب)
عن الألومنيوم (الراشح).
3.~هيدروكسيد المغنيسيوم غير مذبذب: ففي الهيدروكسيد المفرط يبقى
صلبًا، مع هيدروكسيد الحديد(III)، فلا يُفصل شيء.
\end{solution}

\begin{solution}{pb:b1:precipitation:1}
\textbf{1.} \ce{Fe(OH)3(s) <=> Fe^3+ + 3OH-}، $K_s = [\ce{Fe^3+}][\ce{OH-}]^3$؛
\ce{Zn(OH)2(s) <=> Zn^2+ + 2OH-}، $K_s = [\ce{Zn^2+}][\ce{OH-}]^2$؛
\ce{Mg(OH)2(s) <=> Mg^2+ + 2OH-}، وبالمثل.
\textbf{2.} ليست للثوابت الصيغة نفسها (القوة 3 أو 2 للتركيز
$[\ce{OH-}]$) والأيونات الثلاثة بتراكيز مختلفة؛ فلا يمكن مقارنة إلا
قيم pH البدء.
\textbf{3.} عند البدء $c[\ce{OH-}]^n = K_s$، فيكون $n\log[\ce{OH-}] =
-\mathrm{p}K_s - \log c$ و$\mathrm{pH} = \mathrm{p}K_e +
\log[\ce{OH-}] = \mathrm{p}K_e - \frac1n(\mathrm{p}K_s + \log c)$.
\textbf{4.} $14.00 - \frac13(38.55 - 3.00) = 2.15$.
\textbf{5.} $14.00 - \frac12(16.38 - 2.30) = 6.96$.
\textbf{6.} $14.00 - \frac12(11.25 - 1.70) = 9.22$.
\textbf{7.} الحديد(III) (2.15)، ثم الزنك (6.96)، ثم المغنيسيوم (9.22). وعند
pH 2.0، $Q_s = 10^{-3} \times 10^{-36} = 10^{-39} < 10^{-38.55}$ للحديد،
والأيونان الآخران أبعد عن الترسيب: لا شيء صلب.
\textbf{8.} \qty{1.0e-6}{mol/L}.
\textbf{9.} $14.00 - \frac13(38.55 - 6.00) = 3.15$.
\textbf{10.} $\log[\ce{Fe^3+}] = -38.55 + 3(14.00 - \mathrm{pH}) = 3.45 -
3\,\mathrm{pH}$: الميل $-3$.
\textbf{11.} حتى بدء ترسيب هيدروكسيد الزنك، أي pH 6.96.
\textbf{12.} من 3.15 إلى 6.96.
\textbf{13.} \ce{Fe^3+ + OH- <=> FeOH^2+}، $\beta_1 =
[\ce{FeOH^2+}]/([\ce{Fe^3+}][\ce{OH-}]) = 10^{11.82}$.
\textbf{14.} يتساوى النوعان حين $[\ce{OH-}] = 10^{-11.82}$، أي عند pH 2.18:
يغلب \ce{Fe^3+} تحته، و\ce{FeOH^2+} فوقه.
\textbf{15.} $10^{11.82 + 3.15 - 14.00} = 10^{0.97} = 9.3$: عند pH 3.15 يكون
الحديد المذاب 10.3 أضعاف \ce{Fe^3+} الحر، أي \qty{1.0e-5}{mol/L}:
فلا يُزال إلا \qty{99.0}{\percent}. لم يكن التقدير الأول مأمونًا.
\textbf{16.} $[\ce{Fe}]_{\mathrm{dissolved}} = 10^{3.45 - 3\,\mathrm{pH}}\,
(1 + 10^{\mathrm{pH} - 2.18})$.
\textbf{17.} $\log[\ce{FeOH^2+}] = 11.82 + \log[\ce{Fe^3+}] + \mathrm{pH} -
14.00 = 1.27 - 2\,\mathrm{pH}$: الميل $-2$.
\textbf{18.} $1.27 - 2\,\mathrm{pH} \approx -6.00$ يعطي 3.64؛ ومع إدخال
حد \ce{Fe^3+}، $10^{3.45 - 10.92}(1 + 10^{1.46}) =
\num{1.0e-6}$ عند $\mathrm{pH} = 3.64$.
\textbf{19.} $[\ce{Zn^2+}] = \num{5.0e-5}$: $14.00 - \frac12(16.38 - 4.30) =
7.96$. ويبدأ هيدروكسيد المغنيسيوم عند 9.22: فلا.
\textbf{20.} \ce{Zn(OH)2(s) + 2OH- <=> Zn(OH)4^2-}، $K = \beta_4 K_s =
10^{14.45 - 16.38} = 10^{-1.93}$.
\textbf{21.} $[\ce{OH-}]^2 = \num{5.0e-3}/10^{-1.93}$، $\mathrm{pH} =
13.81$. فيجب ألا يتجاوز المشغّل الحد بالقاعدة: فوق pH 9.2 يترسب
المغنيسيوم مع الزنك، وفوق ذلك بكثير يعود الزنك
إلى المحلول.
\textbf{22.} $\num{1.0e-3} \times 1000 \times 0.999 \times 106.8 =
\qty{107}{g}$ لكل متر مكعب.
\textbf{23.} $\num{5.0e-3} \times 1000 \times 0.99 \times 99.4 =
\qty{492}{g}$ لكل متر مكعب.
\textbf{24.} إذا رُفع pH دفعة واحدة فوق 6.96، ترسب الهيدروكسيدان في
حمأة واحدة: فيضيع الزنك في نفايات الحديد، وتتلوث النفايات
بالزنك. أما التوقف بين العتبتين فيزيل الحديد
وحده.
\textbf{25.} \textbf{من pH 3.64 إلى pH 6.96}: تحت 3.64 يبقى أكثر من
\qty{0.1}{\percent} من الحديد مذابًا، على هيئة \ce{FeOH^2+} في الغالب؛ وفوق
6.96 يتكوّن هيدروكسيد الزنك. ولو أُهمل \ce{FeOH^2+} لوُضع
الحد الأدنى عند 3.15.
\end{solution}
